\documentclass[UTF8,11pt,a4paper,oneside,fontset=none]{ctexart}

\usepackage[a4paper,top=1.72cm,bottom=1.72cm,left=1.82cm,right=1.82cm,headheight=15pt]{geometry}
\usepackage{fontspec}
\usepackage{amsmath,amssymb,mathtools,bm}
\usepackage{unicode-math}
\usepackage{booktabs,tabularx,array}
\usepackage[table]{xcolor}
\usepackage{graphicx,float}
\usepackage{siunitx}
\usepackage{enumitem}
\usepackage{fancyhdr}
\usepackage{caption}
\usepackage{hyperref}
\usepackage[most]{tcolorbox}
\usepackage{tikz}
\usetikzlibrary{arrows.meta,positioning}

\setmainfont{Times New Roman}
\setsansfont{Arial}
\setmonofont{Consolas}
\setmathfont{Cambria Math}
\setCJKmainfont[BoldFont=SimHei,ItalicFont=KaiTi]{SimSun}
\setCJKsansfont[BoldFont={Microsoft YaHei Bold}]{Microsoft YaHei}
\setCJKmonofont{FangSong}

\definecolor{DeepBlue}{HTML}{17324D}
\definecolor{PandaTeal}{HTML}{0B7A75}
\definecolor{WarmGold}{HTML}{C58A19}
\definecolor{Coral}{HTML}{C95F4A}
\definecolor{SoftBlue}{HTML}{EEF5F8}
\definecolor{SoftTeal}{HTML}{ECF7F5}
\definecolor{SoftGold}{HTML}{FFF7E7}
\definecolor{SoftRed}{HTML}{FFF0F0}
\definecolor{SoftGray}{HTML}{F5F7F9}
\definecolor{MidGray}{HTML}{667085}
\definecolor{LineGray}{HTML}{D5DEE5}

\hypersetup{
  colorlinks=true,
  linkcolor=DeepBlue,
  urlcolor=PandaTeal,
  pdfauthor={范子恒},
  pdftitle={PandaX-4T run10972 数据结果分析报告}
}

\ctexset{
  section={
    format=\Large\bfseries\color{DeepBlue},
    number=\arabic{section},
    beforeskip=0pt,
    afterskip=0.62em
  },
  subsection={format=\large\bfseries\color{PandaTeal}}
}

\pagestyle{fancy}
\fancyhf{}
\fancyhead[L]{\small\color{DeepBlue}PandaX-4T run10972}
\fancyhead[R]{\small\color{MidGray}数据结果分析报告}
\fancyfoot[C]{\small\color{MidGray}\thepage}
\renewcommand{\headrulewidth}{0.4pt}
\renewcommand{\footrulewidth}{0pt}

\captionsetup{font=small,labelfont=bf,labelsep=quad}
\setlength{\parindent}{2em}
\setlength{\parskip}{0.15em}
\setlist[itemize]{leftmargin=1.55em,itemsep=0.20em,topsep=0.20em,parsep=0pt}
\setlist[enumerate]{leftmargin=1.8em,itemsep=0.22em,topsep=0.20em,parsep=0pt}
\renewcommand{\arraystretch}{1.16}
\sisetup{group-separator={,},group-minimum-digits=4}

\newcolumntype{Y}{>{\raggedright\arraybackslash}X}
\newcolumntype{C}[1]{>{\centering\arraybackslash}p{#1}}
\newcolumntype{R}[1]{>{\raggedleft\arraybackslash}p{#1}}

\newtcolorbox{infobox}[1][]{
  enhanced,breakable,colback=SoftBlue,colframe=DeepBlue,
  boxrule=0.7pt,arc=1.7mm,left=1.5mm,right=1.5mm,top=1.0mm,bottom=1.0mm,#1
}
\newtcolorbox{resultbox}[1][]{
  enhanced,breakable,colback=SoftTeal,colframe=PandaTeal,
  boxrule=0.8pt,arc=1.7mm,left=1.7mm,right=1.7mm,top=1.1mm,bottom=1.1mm,#1
}
\newtcolorbox{warningbox}[1][]{
  enhanced,breakable,colback=SoftGold,colframe=WarmGold,
  boxrule=0.75pt,arc=1.7mm,left=1.5mm,right=1.5mm,top=1.0mm,bottom=1.0mm,#1
}
\newtcolorbox{boundarybox}[1][]{
  enhanced,breakable,colback=SoftRed,colframe=Coral,
  boxrule=0.75pt,arc=1.7mm,left=1.5mm,right=1.5mm,top=1.0mm,bottom=1.0mm,#1
}
\newtcolorbox{graybox}[1][]{
  enhanced,breakable,colback=SoftGray,colframe=LineGray,
  boxrule=0.65pt,arc=1.4mm,left=1.5mm,right=1.5mm,top=1.0mm,bottom=1.0mm,#1
}

\newcommand{\Rsixtyeight}{R_{68}}
\newcommand{\Rninty}{R_{90}}
\newcommand{\OOF}{\mathrm{OOF}}
\newcommand{\good}{\cellcolor{SoftTeal}}
\newcommand{\bad}{\cellcolor{SoftRed}}

\begin{document}

% =============================================================================
% 第 1 页：封面
% =============================================================================
\thispagestyle{empty}
\vspace*{0.68cm}
\noindent{\color{PandaTeal}\rule{\textwidth}{3.5pt}}

\vspace{0.80cm}
{\sffamily\bfseries\fontsize{25pt}{32pt}\selectfont\color{DeepBlue}
PandaX-4T run10972\par
数据结果分析报告\par}

\vspace{0.28cm}
{\sffamily\bfseries\fontsize{15.5pt}{21pt}\selectfont\color{PandaTeal}
事件级五折 OOF 重建结果、稳定性与证据边界\par}

\vspace{0.68cm}
\begin{tcolorbox}[
  colback=SoftBlue,colframe=DeepBlue,boxrule=0.9pt,arc=2mm,
  left=3mm,right=3mm,top=2.5mm,bottom=2.5mm
]
\centering
{\sffamily\bfseries\large 本报告集中回答三个问题}\par
\vspace{1mm}
重建后能量分布改善了多少？\par
这种改善是否在五个独立外折上重复出现？\par
现有证据能够支持哪些结论，又不能支持哪些结论？
\end{tcolorbox}

\vspace{0.60cm}
\begin{center}
\begin{tabular}{C{4.65cm} C{4.65cm} C{4.65cm}}
\begin{tcolorbox}[width=4.4cm,colback=white,colframe=LineGray,arc=1.5mm,halign=center]
{\sffamily\bfseries\fontsize{22pt}{24pt}\selectfont\color{DeepBlue}4500}\par
\small 事件进入 OOF 评价
\end{tcolorbox}
&
\begin{tcolorbox}[width=4.4cm,colback=white,colframe=LineGray,arc=1.5mm,halign=center]
{\sffamily\bfseries\fontsize{22pt}{24pt}\selectfont\color{PandaTeal}2}\par
\small 条独立重建管线
\end{tcolorbox}
&
\begin{tcolorbox}[width=4.4cm,colback=white,colframe=LineGray,arc=1.5mm,halign=center]
{\sffamily\bfseries\fontsize{22pt}{24pt}\selectfont\color{WarmGold}5}\par
\small 个事件级外折
\end{tcolorbox}
\end{tabular}
\end{center}

\vspace{0.32cm}
\begin{boundarybox}[title={与前两册的分工},fonttitle=\bfseries]
《原始数据报告》说明数据内容，《数据处理方法报告》说明处理步骤。本册不重复变量清单和代码流程，只呈现\textbf{事件级结果、逐折稳定性、统计区间、拟合审计与结论边界}。
\end{boundarybox}

\vfill
\noindent
\begin{minipage}[t]{0.48\textwidth}
\textbf{报告人：}范子恒\\
\textbf{数据：}run10972
\end{minipage}\hfill
\begin{minipage}[t]{0.48\textwidth}
\raggedleft
\textbf{日期：}2026 年 7 月 18 日\\
\textbf{版本：}事件级结果版
\end{minipage}

\vspace{0.32cm}
\noindent{\color{PandaTeal}\rule{\textwidth}{1.3pt}}
\clearpage

% =============================================================================
% 第 2 页：结果仪表盘
% =============================================================================
\section{先看结论：结果仪表盘}

\begin{resultbox}[title={两个核心结果},fonttitle=\bfseries]
\begin{center}
\begin{tabular}{C{5.0cm} C{4.2cm} C{4.2cm}}
\toprule
 & \(\Rsixtyeight\) & \(\Rninty\)\\
\midrule
\textbf{A4 联合通道校正} &
\(\mathbf{7.842\%\rightarrow6.748\%}\) &
\(\mathbf{14.723\%\rightarrow13.340\%}\)\\
\textbf{basic 嵌套优胜方案} &
\(\mathbf{8.025\%\rightarrow5.874\%}\) &
\(\mathbf{15.063\%\rightarrow11.899\%}\)\\
\bottomrule
\end{tabular}
\end{center}
\end{resultbox}

\begin{table}[H]
\centering
\small
\caption{固定随机种子的事件级五折 OOF 汇总。改善量定义为校正后减校正前，负值表示分布缩窄。}
\begin{tabular}{@{}C{3.7cm} C{2.15cm} C{2.15cm} C{2.25cm} C{2.25cm} C{2.0cm}@{}}
\toprule
管线 & 原始值 & 校正值 & 改善量 & 相对缩窄 & 五折结论\\
\midrule
A4：\(\Rsixtyeight\) & 7.8422\% & \good 6.7480\% & \good $-1.0941$ pp & \good 13.95\% & 5/5 改善\\
A4：\(\Rninty\) & 14.7232\% & \good 13.3401\% & \good $-1.3830$ pp & \good 9.39\% & 5/5 改善\\
\addlinespace
basic：\(\Rsixtyeight\) & 8.0253\% & \good 5.8738\% & \good $-2.1515$ pp & \good 26.81\% & 5/5 改善\\
basic：\(\Rninty\) & 15.0635\% & \good 11.8989\% & \good $-3.1646$ pp & \good 21.01\% & 5/5 改善\\
\bottomrule
\end{tabular}
\end{table}

\begin{warningbox}[title={两个 raw 基线不能混算},fonttitle=\bfseries]
两条管线的 raw 数值略有差异，原因是有效事件处理规则、归一化方式与权重定义链并不完全相同。因此，本报告只做同一管线内部的前后比较：
\[
\text{A4 对比 A4 的 raw},\qquad
\text{basic 优胜方案对比 basic 的 raw}.
\]
不能用两条校正后曲线的差直接声称某条管线净胜另一条管线。
\end{warningbox}

\begin{infobox}[title={阶段性判断},fonttitle=\bfseries]
\begin{itemize}
  \item 两条管线均同时改善中心宽度 \(\Rsixtyeight\) 与更宽尾部尺度 \(\Rninty\)，不是只把峰中心局部压窄。
  \item 两条管线在五个外折中均逐折改善两个主指标，结果不是由单一外折驱动。
  \item A4 结构简单、物理含义清楚；basic 优胜方案数值改善更大，但模型自由度也更高。
\end{itemize}
\end{infobox}
\clearpage

% =============================================================================
% 第 3 页：指标和协议
% =============================================================================
\section{评价口径与事件级 OOF 协议}

\subsection*{2.1 为什么同时报告 \(\Rsixtyeight\) 和 \(\Rninty\)}

对拼接后的全部外折测试事件，以相对能量分位数 \(Q_p\) 定义
\begin{align}
\Rsixtyeight
  &=\frac{Q_{84.135}-Q_{15.865}}{2Q_{50}},\\
\Rninty
  &=\frac{Q_{95}-Q_{5}}{2Q_{50}}.
\end{align}
\(\Rsixtyeight\) 主要描述分布中央约 68\% 的宽度，\(\Rninty\) 对较远尾部更敏感。二者都不要求能谱服从高斯分布，因此适合作为当前探索阶段的主指标。

改善量及相对缩窄比例分别为
\begin{align}
\Delta R &= R_{\mathrm{corrected}}-R_{\mathrm{raw}},\\
I_R &= \frac{R_{\mathrm{raw}}-R_{\mathrm{corrected}}}{R_{\mathrm{raw}}}\times100\%.
\end{align}
这里 \(\Delta R<0\) 表示改善；pp 表示百分点，而不是相对百分比。

\subsection*{2.2 外折测试事件从不参与训练}

\begin{center}

\begin{tikzpicture}[
  node distance=9mm,
  box/.style={draw=DeepBlue,rounded corners=1.5mm,fill=SoftBlue,
              minimum width=4.15cm,minimum height=1.25cm,align=center,font=\small},
  arrow/.style={-{Latex[length=2.2mm]},thick,color=PandaTeal}
]
\node[box] (split) {固定种子、均衡的\\事件级五折划分};
\node[box,right=of split] (train) {仅在四折内\\学习响应与权重};
\node[box,right=of train] (test) {在留出折产生一次\\OOF 能量预测};
\draw[arrow] (split) -- (train);
\draw[arrow] (train) -- (test);
\end{tikzpicture}
\end{center}

4500 个事件被均衡分到五折，每折 900 个事件。每一轮只用四折训练，在剩余一折上预测；五轮结束后，每个事件恰好得到一次未参与自身训练的 OOF 预测。所有主指标均在 4500 个 OOF 预测拼接后计算。

\begin{graybox}[title={如何理解“稳定”},fonttitle=\bfseries]
本报告要求三个层次同时成立：
\begin{enumerate}
  \item 全部 OOF 事件拼接后的 \(\Rsixtyeight\) 与 \(\Rninty\) 均改善；
  \item 五个外折中的两个指标均逐折改善；
  \item 配对事件重采样得到的改善区间不跨越零。
\end{enumerate}
这比只展示一张校正后能谱更严格，但仍然只是同一 run 内的事件级证据。
\end{graybox}

\begin{boundarybox}[title={本协议没有证明什么},fonttitle=\bfseries]
事件级 OOF 检验不能替代独立 run。它不能单独证明跨运行条件、跨时间段或跨能区泛化；也不能消除上游样本选择对能谱形状的潜在影响。
\end{boundarybox}
\clearpage

% =============================================================================
% 第 4 页：channel 消融
% =============================================================================
\section{Channel physics 管线：消融结果}

该管线先分别处理 S1 与底部 S2 的低自由度响应，再以正权重组合相对能量。A0--A4 的作用是逐步回答：单独修正一个通道是否足够，以及两个通道是否必须联合处理。

\begin{figure}[H]
  \centering
  \includegraphics[width=0.82\textwidth]{../channel_physics_outputs/channel_physics_matrix.png}
  \caption{Channel physics 候选的事件级 OOF 指标。数值越低表示全事件分布越窄。}
\end{figure}

\begin{table}[H]
\centering
\small
\caption{Channel physics 消融矩阵。A0 是本管线 raw 基线。}
\begin{tabularx}{\textwidth}{@{}C{1.25cm}Y C{2.25cm} C{2.25cm} C{2.25cm}@{}}
\toprule
候选 & 处理内容 & \(\Rsixtyeight\) & \(\Rninty\) & 相对 A0 判断\\
\midrule
A0 & 原始 S1 与原始底部 S2 的正权组合 & 7.8422\% & 14.7232\% & 基线\\
A1 & 只处理底部 S2 的漂移时间响应 & \bad 9.3895\% & \bad 17.0109\% & 两项均变差\\
A2 & 只处理 S1 的漂移时间与半径响应 & \bad 9.0931\% & \bad 17.3172\% & 两项均变差\\
A3 & 只处理底部 S2 的漂移时间与半径响应 & \bad 8.5892\% & \bad 15.6796\% & 两项均变差\\
A4 & 同时处理 S1 与底部 S2 的响应 & \good 6.7480\% & \good 13.3401\% & 两项均改善\\
\bottomrule
\end{tabularx}
\end{table}

\begin{resultbox}[title={消融实验真正说明了什么？},fonttitle=\bfseries]
\begin{itemize}
  \item A1--A3 的负结果表明：只改一个通道会破坏原始 S1--S2 之间已有的互补关系，不能把“单通道响应更平坦”直接等同于“组合能量更窄”。
  \item A4 同时修正两个通道后才出现稳定增益，说明关键不只是某一个变量，而是\textbf{两条通道响应需要成套处理}。
  \item 因此，A4 的改进不是从消融矩阵中挑一个偶然最小值，而是由清楚的正负对照支持。
\end{itemize}
\end{resultbox}

\begin{warningbox}[title={不要删除负结果},fonttitle=\bfseries]
A1--A3 虽然没有提高分辨率，却限制了后续解释：下一步不应继续沿着“只加强某一单通道修正”的方向堆叠复杂度，而应检查联合响应中仍残留的二维结构。
\end{warningbox}
\clearpage

% =============================================================================
% 第 5 页：A4 稳定性
% =============================================================================
\section{A4：逐折稳定性与条件区间}

\begin{table}[H]
\centering
\small
\caption{A4 五个外折的测试侧结果。所有数值均为百分数，\(\Delta\) 为校正后减校正前。}
\begin{tabular}{@{}C{1.0cm} C{2.0cm} C{2.0cm} C{1.9cm} C{2.0cm} C{2.0cm} C{1.9cm}@{}}
\toprule
折 & raw \(\Rsixtyeight\) & A4 \(\Rsixtyeight\) & \(\Delta\Rsixtyeight\) & raw \(\Rninty\) & A4 \(\Rninty\) & \(\Delta\Rninty\)\\
\midrule
0 & 8.030 & \good 6.724 & \good $-1.306$ & 14.730 & \good 13.714 & \good $-1.015$\\
1 & 7.595 & \good 6.296 & \good $-1.298$ & 14.297 & \good 12.597 & \good $-1.699$\\
2 & 8.497 & \good 8.049 & \good $-0.448$ & 15.175 & \good 14.778 & \good $-0.397$\\
3 & 7.334 & \good 6.693 & \good $-0.641$ & 14.632 & \good 12.688 & \good $-1.944$\\
4 & 7.921 & \good 6.758 & \good $-1.163$ & 14.146 & \good 13.331 & \good $-0.816$\\
\bottomrule
\end{tabular}
\end{table}

五折均满足 \(\Delta\Rsixtyeight<0\) 且 \(\Delta\Rninty<0\)。第 2 折的改善最小，但仍保持同方向；这说明总体结果并非由少数特别容易的测试事件主导，同时也提示不同事件子集上的收益并不完全相同。

\begin{resultbox}[title={A4 的配对事件 bootstrap},fonttitle=\bfseries]
在冻结的 OOF 预测上进行 1000 次配对事件重采样：
\begin{align*}
\operatorname{median}(\Delta\Rsixtyeight)
  &= -1.0991\ \text{pp}, &
95\%\ \mathrm{CI}&=[-1.3034,-0.8955]\ \text{pp},\\
\operatorname{median}(\Delta\Rninty)
  &= -1.4199\ \text{pp}, &
95\%\ \mathrm{CI}&=[-1.8062,-1.0552]\ \text{pp}.
\end{align*}
两个区间都完全位于零以下，支持“改善不是有限事件抽样下的随机符号波动”。
\end{resultbox}

\subsection*{4.1 响应平坦度与尾部的补充检查}

\begin{table}[H]
\centering
\small
\caption{A4 前后的 OOF 诊断。峰峰变化越小表示对应方向的剩余响应趋势越弱。}
\begin{tabular}{@{}C{5.8cm} C{3.2cm} C{3.2cm} C{2.2cm}@{}}
\toprule
诊断量 & raw & A4 & 变化\\
\midrule
漂移时间方向峰峰变化 & 3.973\% & \good 2.379\% & 降低\\
半径平方方向峰峰变化 & 10.230\% & \good 5.230\% & 约减半\\
固定 $\pm20\%$ 范围外比例 & 3.733\% & \good 2.556\% & 降低\\
相对 $2\Rsixtyeight$ 之外比例 & 8.156\% & 9.622\% & 上升\\
\bottomrule
\end{tabular}
\end{table}

最后一行并不与前面矛盾：A4 使 \(\Rsixtyeight\) 本身变小，因而“\(2\Rsixtyeight\)”对应的动态阈值也随之收紧。固定范围尾部减少，而相对收紧后核心宽度定义的尾部比例上升，说明主峰变窄后仍存在非高斯尾部。
\clearpage

% =============================================================================
% 第 6 页：basic 方法矩阵
% =============================================================================
\section{Basic matrix：方法搜索结果}

basic 管线在外层训练数据内部选择正的 S1--S2 权重与低阶响应形式，再在留出事件上产生 OOF 预测。嵌套优胜方案在五折中均选择 \texttt{dt\_r2\_then\_xy}。

\begin{figure}[H]
  \centering
  \includegraphics[width=0.84\textwidth]{../basic_matrix_outputs/method_matrix.png}
  \caption{Basic 方法矩阵的事件级 OOF 比较。红色优胜方案由外层训练数据内部的选择规则产生。}
\end{figure}

\begin{table}[H]
\centering
\scriptsize
\renewcommand{\arraystretch}{1.08}
\caption{Basic 管线全部候选的 OOF 主指标，按 \(\Rsixtyeight\) 从小到大排列。}
\begin{tabular}{@{}C{4.25cm} C{2.6cm} C{2.6cm} C{4.25cm}@{}}
\toprule
候选 & \(\Rsixtyeight\) & \(\Rninty\) & 结果解释\\
\midrule
\texttt{dt\_r2\_then\_xy} & \good 5.8738\% & \good 11.8989\% & 嵌套优胜方案\\
\texttt{r2\_then\_dt} & 6.2100\% & 12.1892\% & 次优\\
\texttt{dt\_r2\_joint} & 6.2889\% & 12.6247\% & 改善\\
\texttt{r2\_quad} & 6.3491\% & 12.4357\% & 改善\\
\texttt{dt\_then\_r2} & 6.3647\% & 12.4515\% & 改善\\
\texttt{dt\_then\_xy} & 6.6055\% & 12.9397\% & 改善\\
\texttt{xy\_quad} & 6.6886\% & 13.2191\% & 改善\\
\texttt{dt\_quad} & 7.5626\% & 14.3830\% & 小幅改善\\
raw & 8.0253\% & 15.0635\% & 本管线基线\\
\bottomrule
\end{tabular}
\end{table}

\begin{resultbox}[title={basic 优胜方案的全样本 OOF 改善},fonttitle=\bfseries]
\begin{align*}
\Rsixtyeight &: 8.0253\%\rightarrow5.8738\%,
&\Delta\Rsixtyeight&=-2.1515\ \text{pp},
&I_{68}&=26.81\%,\\
\Rninty &: 15.0635\%\rightarrow11.8989\%,
&\Delta\Rninty&=-3.1646\ \text{pp},
&I_{90}&=21.01\%.
\end{align*}
数值收益明显大于 A4，但这种差异不能直接解释为净方法优势，因为两条管线的 raw 定义链并不完全相同。
\end{resultbox}
\clearpage

% =============================================================================
% 第 7 页：basic 稳定性
% =============================================================================
\section{Basic 优胜方案：逐折稳定性}

\begin{figure}[H]
  \centering
  \includegraphics[width=0.62\textwidth]{../basic_matrix_outputs/fold_stability.png}
  \caption{Basic 候选在五个事件级外折上的 \(\Rsixtyeight\)。该图用于检查改进是否集中在个别外折。}
\end{figure}

\begin{table}[H]
\centering
\scriptsize
\renewcommand{\arraystretch}{1.10}
\caption{Basic raw 与嵌套优胜方案的逐折结果。所有外折均选择 \texttt{dt\_r2\_then\_xy}。}
\begin{tabular}{@{}C{0.8cm} C{1.7cm} C{1.7cm} C{1.7cm} C{1.7cm} C{1.7cm} C{1.7cm} C{1.7cm}@{}}
\toprule
折 & 选择的 $\alpha$ & raw \(\Rsixtyeight\) & 优胜 \(\Rsixtyeight\) & \(\Delta\Rsixtyeight\) & raw \(\Rninty\) & 优胜 \(\Rninty\) & \(\Delta\Rninty\)\\
\midrule
0 & 0.652 & 8.312 & \good 6.275 & \good $-2.037$ & 15.119 & \good 11.556 & \good $-3.564$\\
1 & 0.425 & 7.672 & \good 5.390 & \good $-2.282$ & 14.607 & \good 11.333 & \good $-3.274$\\
2 & 0.652 & 8.541 & \good 6.528 & \good $-2.013$ & 15.177 & \good 12.289 & \good $-2.887$\\
3 & 0.425 & 7.730 & \good 5.889 & \good $-1.840$ & 14.683 & \good 11.803 & \good $-2.880$\\
4 & 0.425 & 7.691 & \good 5.502 & \good $-2.188$ & 14.795 & \good 11.407 & \good $-3.388$\\
\bottomrule
\end{tabular}
\end{table}

\begin{resultbox}[title={Basic 优胜方案的配对事件 bootstrap},fonttitle=\bfseries]
\begin{align*}
\operatorname{median}(\Delta\Rsixtyeight)
  &= -2.1477\ \text{pp}, &
95\%\ \mathrm{CI}&=[-2.4558,-1.8803]\ \text{pp},\\
\operatorname{median}(\Delta\Rninty)
  &= -3.1732\ \text{pp}, &
95\%\ \mathrm{CI}&=[-3.6608,-2.6607]\ \text{pp}.
\end{align*}
两个区间均不跨零，且每折改善方向一致。就当前 run 内的事件级检验而言，数值改善具有较强的一致性。
\end{resultbox}

\begin{warningbox}[title={方法选择仍然需要独立验证},fonttitle=\bfseries]
五折都选择同一种响应形式，说明内层选择相当一致；但候选族是在当前研究阶段建立并检查的。要把它升级为正式默认方案，仍需先冻结全部规则，再在独立 run 上一次性验证。
\end{warningbox}
\clearpage

% =============================================================================
% 第 8 页：bootstrap 和 GOF
% =============================================================================
\section{统计稳定性与峰形拟合审计}

\subsection*{7.1 两条管线的事件重采样结果}

\begin{table}[H]
\centering
\small
\caption{冻结 OOF 预测上的配对事件 bootstrap。区间全部位于零以下。}
\begin{tabular}{@{}C{3.1cm} C{2.2cm} C{2.6cm} C{3.0cm} C{3.0cm}@{}}
\toprule
管线 & 指标 & 中位改善量 & 95\% 区间下限 & 95\% 区间上限\\
\midrule
A4 & \(\Delta\Rsixtyeight\) & $-1.0991$ pp & $-1.3034$ pp & $-0.8955$ pp\\
A4 & \(\Delta\Rninty\) & $-1.4199$ pp & $-1.8062$ pp & $-1.0552$ pp\\
basic & \(\Delta\Rsixtyeight\) & $-2.1477$ pp & $-2.4558$ pp & $-1.8803$ pp\\
basic & \(\Delta\Rninty\) & $-3.1732$ pp & $-3.6608$ pp & $-2.6607$ pp\\
\bottomrule
\end{tabular}
\end{table}

\begin{warningbox}[title={bootstrap 区间的准确含义},fonttitle=\bfseries]
每次重采样都对同一批 raw 与 corrected 事件使用完全相同的抽样索引，因此保留了逐事件配对关系。但这里没有在每次复制中重新训练响应图、权重和方法选择；区间是\textbf{冻结 OOF 预测的条件区间}。此外，事件级 bootstrap 假设事件可以交换，不能表达 run 间系统变化。
\end{warningbox}

\subsection*{7.2 为什么高斯 core 不是正式结果}

作为补充诊断，在固定相对能量窗口内尝试“高斯信号加二阶 Bernstein 背景”，并用 200 次参数 bootstrap 检验拟合优度。

\begin{table}[H]
\centering
\small
\caption{校正后 OOF 能谱的峰形拟合审计。}
\begin{tabular}{@{}C{3.5cm} C{3.4cm} C{3.0cm} C{4.2cm}@{}}
\toprule
管线 & core \(\sigma/\mu\) & 拟合优度 $p$ 值 & 判定\\
\midrule
A4 & 3.6543\% & 0.004975 & \bad 拒绝该峰形模型\\
basic 优胜方案 & 4.1548\% & 0.0199005 & \bad 拒绝该峰形模型\\
\bottomrule
\end{tabular}
\end{table}

两项 $p$ 值都小于 0.05，说明该参数化模型不能充分描述真实 OOF 分布。于是
\[
\boxed{\text{正式结论继续采用全事件非参数 }\Rsixtyeight\text{ 与 }\Rninty}
\]
而不能把 3.6543\% 或 4.1548\% 当作已经成立的能量分辨率。basic 的 \(\Rsixtyeight\) 更小但拟合 core 更大，也恰好说明：当峰形含有背景和非高斯结构时，不同宽度定义不能互相替代。

\begin{infobox}[title={这项“拒绝”是有价值的结果},fonttitle=\bfseries]
拟合失败并不否定 OOF 分布确实缩窄；它只否定“用当前高斯加背景形式概括全部峰形”的假设。保留这一负结果可以避免报告一个看似更漂亮、却没有通过拟合优度检验的数字。
\end{infobox}
\clearpage

% =============================================================================
% 第 9 页：综合解释
% =============================================================================
\section{综合解释：我们已经知道什么}

\begin{resultbox}[title={证据层级一：直接观察成立},fonttitle=\bfseries]
\begin{itemize}
  \item A4 将 \(\Rsixtyeight\) 从 7.8422\% 降到 6.7480\%，将 \(\Rninty\) 从 14.7232\% 降到 13.3401\%。
  \item basic 优胜方案将自身 \(\Rsixtyeight\) 从 8.0253\% 降到 5.8738\%，将 \(\Rninty\) 从 15.0635\% 降到 11.8989\%。
  \item A4 的消融实验显示，单独处理任一通道均不足，联合处理才产生增益。
\end{itemize}
\end{resultbox}

\begin{infobox}[title={证据层级二：同一 run 内稳定},fonttitle=\bfseries]
\begin{itemize}
  \item 两条管线的五个外折均同时改善 \(\Rsixtyeight\) 和 \(\Rninty\)。
  \item basic 的五个外折均选择 \texttt{dt\_r2\_then\_xy}，方法选择没有在外折之间来回跳变。
  \item 四个 bootstrap 改善区间都不跨零，支持改善方向的事件级稳定性。
\end{itemize}
\end{infobox}

\begin{boundarybox}[title={证据层级三：尚未证明跨 run 泛化},fonttitle=\bfseries]
当前结果仍来自单一 run 的 4500 个上游标量事件。事件级 OOF 能阻止同一事件参与自身训练，却不能模拟新的运行条件。输入名称还暗示样本可能经历了上游单散射与能量阈值选择，因此现有数字只描述给定样本，不能直接外推到全部触发事件或其他能区。
\end{boundarybox}

\subsection*{8.1 两条管线应怎样定位}

\begin{table}[H]
\centering
\small
\caption{当前阶段的功能定位，而不是未经统一基线的胜负排名。}
\begin{tabularx}{\textwidth}{@{}C{3.0cm}Y Y@{}}
\toprule
管线 & 优点 & 当前限制\\
\midrule
A4 & 通道级物理结构清楚；消融关系明确；自由度较低；便于解释剩余响应 & 数值改善较温和；仍存在非高斯尾部与位置残余\\
basic 优胜方案 & 当前事件级 OOF 的 \(\Rsixtyeight\) 与 \(\Rninty\) 最低；五折选择一致 & 响应链更复杂；尚未在统一 raw 定义和独立 run 上确认净优势\\
\bottomrule
\end{tabularx}
\end{table}

\begin{warningbox}[title={最稳妥的阶段性表述},fonttitle=\bfseries]
可以说“两个独立构造的低自由度重建方案都在 run10972 的事件级 OOF 检验中稳定缩窄了全事件分布”；暂时不能说“已经获得最终探测器能量分辨率”，也不能说“basic 已被证明普遍优于 A4”。
\end{warningbox}
\clearpage

% =============================================================================
% 第 10 页：结论与下一步
% =============================================================================
\section{最终结论与下一步工作}

\begin{resultbox}[title={本阶段已经做到的事情},fonttitle=\bfseries]
\begin{enumerate}
  \item 从基础 S1、底部 S2、漂移时间和横向位置出发，建立了两条相互独立的相对能量重建管线。
  \item 用固定种子、均衡的事件级五折 OOF 协议评价全部 4500 个事件，避免测试事件参与自身模型训练。
  \item A4 在全样本上使 \(\Rsixtyeight\) 相对缩窄 13.95\%、\(\Rninty\) 相对缩窄 9.39\%；basic 优胜方案对应缩窄 26.81\% 与 21.01\%。
  \item 两条管线均在五个外折上逐折改善两个指标，配对事件 bootstrap 区间也全部位于零以下。
  \item 峰形拟合优度审计拒绝了当前高斯 core 模型，因此主动保留 \(\Rsixtyeight\) 和 \(\Rninty\) 作为正式结果，没有把未通过检验的 core 数值包装成最终分辨率。
\end{enumerate}
\end{resultbox}

\subsection*{9.1 推荐的下一步顺序}

\begin{enumerate}
  \item \textbf{冻结当前方案。} 固定 A4 与 \texttt{dt\_r2\_then\_xy} 的变量、响应形式、权重搜索范围和评价指标，避免看到新数据后继续调规则。
  \item \textbf{建立统一 raw 基线。} 让两条管线共享相同的有效事件规则、归一化和初始权重定义，再做严格的净增益比较。
  \item \textbf{验证独立 run。} 在新的运行数据上只执行一次冻结流程，重点检查 \(\Rsixtyeight\)、\(\Rninty\)、逐折稳定性与尾部比例是否重复改善。
  \item \textbf{补全上游选择信息。} 明确能量阈值的定义、效率以及目标峰信息，判断当前样本是否已被选择条件塑形。
  \item \textbf{升级不确定度。} 在计算资源允许时，对每次重采样完整重训响应与权重，并增加不同随机种子的敏感性检查。
\end{enumerate}

\subsection*{9.2 可复核的结果入口}

\begin{graybox}
\small
\begin{tabularx}{\textwidth}{@{}Y Y@{}}
Channel 汇总与逐折结果 & \path{../channel_physics_outputs/summary.csv}, \path{../channel_physics_outputs/fold_metrics.csv}\\
Channel 条件区间与 GOF & \path{../channel_physics_outputs/paired_event_bootstrap_summary.csv}, \path{../channel_physics_outputs/robust_peak_corrected_gof.json}\\
Basic 汇总与逐折结果 & \path{../basic_matrix_outputs/summary_metrics.csv}, \path{../basic_matrix_outputs/fold_metrics.csv}\\
Basic 条件区间与 GOF & \path{../basic_matrix_outputs/paired_event_bootstrap_summary.csv}, \path{../basic_matrix_outputs/robust_peak_corrected_gof.json}\\
\end{tabularx}
\end{graybox}

\begin{boundarybox}[title={一句话总结},fonttitle=\bfseries]
run10972 的事件级 OOF 证据表明，A4 与 basic 优胜方案都能稳定缩窄当前样本的中心和宽尾分布；但最终结论仍必须由统一基线和独立 run 的冻结验证来完成。
\end{boundarybox}

\end{document}
